\section{Institutional Background and Related Literature}
\label{sec:background}

\subsection{HMGP as an Intergovernmental Grant}
\label{sec:hmgp}

HMGP is authorized by Section 404 of the Stafford Act and becomes available in a state after a major disaster declaration. The amount available for a declaration is set on a sliding scale, up to 15 percent of the first \$2 billion of estimated federal disaster assistance, with higher percentages for states with enhanced mitigation plans \parencite{horn2022crs}. Although the declaration designates particular counties for disaster assistance, the state manages the mitigation funds as the recipient and decides which subapplications to fund \parencite{horn2022crs}. Eligible subapplicants include state agencies, local governments, and certain nonprofits; households benefit only through a community application. HMGP accounted for about 80 percent of FEMA's Hazard Mitigation Assistance spending between 1989 and 2019 \parencite{frimpong2023federal}.

The program has three features that matter for the distribution of funds across places. The pool of money is set by the declaration, so counties within a declaration draw on the same allocation under the same state administration. Funds are matched, so a subapplicant must find the nonfederal share. And money reaches a place only through a subapplication from an eligible entity located there, so the number and activity of local governments determine how many claims a county can make on the allocation.

For tornado and wind hazards, HMGP finances community and school safe rooms, individual residential safe rooms through state-run programs, and wind retrofits of public and private structures. Tornado fatalities are concentrated among residents of mobile homes and in the Southeast \parencite{ashley2007fatalities,sutter2010mobile}, the cost-effectiveness of shelters depends on housing type and local risk \parencite{merrell2005casualties,simmons2006shelters}, and the geography of tornado activity has been shifting toward the Midwest and Southeast \parencite{gensini2018trends}. FEMA has twice tried to reduce the administrative burden of such projects. Pre-calculated benefits allow residential safe rooms that meet design criteria to forgo a project-specific benefit--cost analysis \parencite{fema2015saferoomprecalc,fema2024bcapolicy}. Section 1215 of the Disaster Recovery Reform Act of 2018 (DRRA) \parencite{drra2018} allows subrecipients to recover management costs of up to 5 percent of their subaward, a provision that FEMA applies to disasters declared on or after August 1, 2017 \parencite{fema1215mgmtcosts}. We use the second reform as a test in Section~\ref{sec:explanations}.

\subsection{Local Government Size, Capacity, and Access to Federal Funds}
\label{sec:literature}

Research on how individuals and organizations respond to hazards emphasizes that action depends on whether a response is judged feasible and within available resources, not only on perceived risk \parencite{lindell2012padm,bubeck2012review}. For a local government, the relevant resource is administrative capacity: staff time, technical expertise, fiscal room for a cost share, and experience with federal programs. Planning research shows that the involvement of professional planners changes the content of local mitigation plans \parencite{lyles2014planners} and that federal and state policies shape local mitigation choices through local commitment and capacity \parencite{burby2006katrina,berke2014impacts}. Local governments learn from their own hazard history, faster under some conditions than others \parencite{brody2009policy}. Organizational capacity predicts local flood mitigation \parencite{brody2010identifying}, and participation in FEMA's Community Rating System rises with local capacity and tax base \parencite{posey2009determinants,sadiq2015local,li2018flood}; interviews with program coordinators and floodplain managers describe why communities do or do not take part \parencite{sadiq2020participation}.

A second body of work studies how federal disaster resources are distributed across places. It builds on evidence that the capacity to prepare for and recover from hazards is unevenly distributed geographically \parencite{cutter2003sovi,cutter2008temporal} and that rural and urban counties differ in measured resilience \parencite{cutter2016urbanrural}. Studies of household assistance report distributive inequities in how aid relates to social vulnerability \parencite{drakes2021social,emrich2022distributive}, and programs that require government applicants tend to favor counties in better socioeconomic condition \parencite{miao2025equity}. Buyouts show disparities by county wealth and race \parencite{mach2019buyouts,elliott2020racial}, and rural HMGP recipients receive less assistance per recipient than urban ones \parencite{seong2022naturally}. In competitive pre-disaster programs, local capacity predicts success \parencite{smith2025rural}, and small and impoverished communities obtained 29 percent of the funding they requested in the first round of FEMA's Building Resilient Infrastructure and Communities program \parencite{headwaters_nd}. The allocation of declarations themselves has a political dimension \parencite{garrett2003political,reeves2011political}, which is one reason we compare counties within a declaration.

States stand between the federal program and local applicants. A national survey of state hazard mitigation officers found large differences in states' ability to help localities develop applications and little use of funding for pre-application work \parencite{smith2020state,vila2022equity}. Research on other grant programs finds that the state's design of the application process can widen or narrow disparities among localities \parencite{collins2006redistributive,collins2008taken}.

This work has not established where in the grant process small jurisdictions fall behind, or how much of any gap reflects the structure of local government rather than the behavior of individual governments. A county with few incorporated places, school districts, and special districts has fewer entities that could apply, so it would obtain fewer projects even if each entity were as likely to apply as its counterparts elsewhere. Separating this structural channel from a constraint on individual governments requires application-level data and a comparison that holds the hazard and the administering state fixed.

\subsection{Expectations}
\label{sec:hypotheses}

We organize the analysis around four expectations. (H1) Conditional on hazard exposure, damage, and the declaration itself, small counties are less likely to obtain HMGP projects. (H2) If the gap reflects local conditions rather than federal or state screening, it arises at the application stage rather than in decisions on submitted requests. (H3) If the gap is not solely a consequence of the number of eligible entities, it persists after accounting for that number and appears for the county government itself. (H4) If local capacity constrains the response to hazard events, a tornado hit raises applications less in small counties than in larger counties, and policies that lower the cost of participating or shift the work to the state are associated with a smaller gap.
